%----------
%	MODELOS GENERATIVOS DE LENGUAJE
%----------

\label{sec:llms_models_iris}

Un modelo generativo de lenguaje, o LLM, es una red
neuronal, habitualmente de tipo \textit{Transformer}~\cite{vaswani2017}, entrenada
sobre grandes volúmenes de texto para predecir la continuación de una secuencia. De ese entrenamiento sale una competencia
lingüística general que puede orientarse a tareas concretas mediante
instrucciones en lenguaje natural, sin entrenar un modelo distinto para cada
tarea. Este trabajo explota justamente esa competencia: la metodología de
codificación se expresa como instrucción y el modelo la aplica a cada pieza.

Para desplegar estos modelos hay dos vías, y la diferencia entre ellas importa en
un proyecto como este.

Los \textbf{modelos locales} son modelos de pesos abiertos que se ejecutan en
infraestructura propia. En este trabajo se usan a través de
\textit{Ollama}\footnote{\url{https://ollama.com} [Último acceso: 05/08/2026]} sobre la granja de GPU del
Departamento de Teoría de la Señal y Comunicaciones. Tienen
dos ventajas principales. La primera es la privacidad: el texto nunca sale de la
infraestructura de la institución, algo relevante cuando se trabaja con contenido
periodístico y datos de investigación. La segunda es el coste: la inferencia no se
paga por uso, a diferencia de las API comerciales, por lo que procesar el corpus
completo no tiene coste marginal. A cambio, exigen hardware especializado y suelen
ser más pequeños y menos capaces que los modelos comerciales punteros.

Los \textbf{modelos comerciales vía API} los alojan las empresas proveedoras (OpenAI,
Google) y se acceden por red. Ofrecen más capacidad y no requieren
infraestructura propia. La contrapartida es doble: se paga por uso y el texto
viaja a un tercero, con lo que ello implica en privacidad.

Comparar ambas vías sobre la misma tarea y el mismo corpus es uno de los
objetivos metodológicos de este trabajo. Para lograrlo, la cadena de procesos de
clasificación es agnóstica a la proveedora: una función de enrutado dirige cada
petición al \textit{backend} que corresponda según el identificador del modelo. De este
modo, lo que se observe al comparar modelos procede de los modelos, no de cómo
está escrito el código que los llama. Los modelos concretos se detallan en la
Sección~\ref{sec:genai_models_iris}.
